\section{Predictive Sufficiency and History-Moment Fusion}
\label{sec:recovery-moments}

Let $S_k$ contain the physical and instrumental state, predictive mode,
and any parameter or dwell-time memory needed to make the dynamics Markov.
Let $H_k$ denote historical cause labels. If
\[
 p(S_{k+1},Y_{k+1}\mid S_{0:k},H_{0:k},Y_{1:k})
 =p(S_{k+1},Y_{k+1}\mid S_k),
\]
then the filtering law $p(S_k\mid Y_{1:k})$ is sufficient for prediction:
for every measurable report set $B$,
\[
 p(Y_{k+1}\in B\mid Y_{1:k})
 =\int p(Y_{k+1}\in B\mid S_k=s)\,p(ds\mid Y_{1:k}).
\]
This is the standard filtering identity. Its condition is substantive: omitted
mode, dwell age, uncertain parameters or instrument memory can make cause
labels predictive even after conditioning on an incomplete state summary.
The marginal state law suffices for forecasting under this factorization;
retrospective cause attribution still requires a joint posterior retaining
$H_k$.

For a finite history mixture with weights $w_j$, let $X$ be a common retained
state vector and $A$ a common scalar severity coordinate. Write
$m_j=E[X\mid j]$, $P_j=\operatorname{Cov}(X\mid j)$,
$\mu_j=E[A\mid j]$, $V_j=\operatorname{Var}(A\mid j)$ and
$C_j=\operatorname{Cov}(X,A\mid j)$. The law of total covariance gives
\begin{align}
 \bar m&=\sum_j w_jm_j, & \bar\mu&=\sum_jw_j\mu_j,\\
 \operatorname{Cov}(X)&=\sum_jw_jP_j+
 \sum_jw_j(m_j-\bar m)(m_j-\bar m)^T,\\
 \operatorname{Cov}(X,A)&=\sum_jw_jC_j+
 \sum_jw_j(m_j-\bar m)(\mu_j-\bar\mu),\\
 \operatorname{Var}(A)&=\sum_jw_jV_j+
 \sum_jw_j(\mu_j-\bar\mu)^2.
\end{align}
Thus a moment-matched Gaussian can retain within-history severity uncertainty
and between-history dispersion without retaining a multimodal mixture. These
identities do not establish that moment matching improves event attribution.

In the controlled study, mean merging preserves the estimated physical mean while using the nominal
conditional covariance; it does not propagate the discarded between-history
and severity covariance. The history-versus-carry contrast therefore cannot
be interpreted as a causal effect of labels alone: it also changes which
recovery alternatives and uncertainty are represented. The study supports its
reported procedure-level comparison, conditional on the tested plants and
protocol.
